
% =============================================================
\section{Physical Foundations}
% =============================================================

The Compton effect describes the scattering of a photon by an initially stationary electron. 
To understand the observed change in wavelength, both the energy 
and the momentum of the photon must be taken into account.

% -------------------------------------------------------------
\subsection{Energy and Momentum of the Photon}
% -------------------------------------------------------------
According to quantum theory, a photon carries the energy
\[
E = h\nu
\]
and the momentum
\[
p = \frac{E}{c} = \frac{h\nu}{c} = \frac{h}{\lambda}.
\]
Although photons have no rest mass, they possess momentum 
that can be transferred during scattering processes. 

\begin{DidacticBox}[Energy in Contrast to the Classical Wave Description]
	In classical wave theory, energy is transported through the amplitude of the wave, 
	not through individual quanta. 
	The Compton effect shows that light occurs in discrete energy packets \( h\nu \) — 
	a clear indication of the particle nature of light.
\end{DidacticBox}

% -------------------------------------------------------------
\subsection{Collision Between Photon and Electron}
% -------------------------------------------------------------
A photon strikes a stationary electron (mass \( m_e \)). 
After the scattering event, the photon has a lower energy \( h\nu' \) 
and is deflected by an angle \( \theta \). 
The electron receives a recoil momentum \( \vec{p}_e \).

\begin{PhysicsBox}[Scattering Geometry]
	\small
	The scattering geometry can be expressed vectorially as
	\[
	\vec{p}_\gamma = \vec{p}_{\gamma'} + \vec{p}_e.
	\]
	In components:
	\[
	p_\gamma = p_{\gamma'}\cos\theta + p_{e,x}, \qquad
	0 = p_{\gamma'}\sin\theta - p_{e,y}.
	\]
\end{PhysicsBox}

\begin{NoteBox}[Inelastic Collision]
	Although energy and momentum are conserved, 
	the collision is \emph{not elastic} in the classical sense, 
	because part of the photon’s energy is converted into the electron’s kinetic energy.
\end{NoteBox}

% -------------------------------------------------------------
\subsection{Energy and Momentum Conservation Combined}
% -------------------------------------------------------------
According to the law of energy conservation:
\[
h\nu + m_e c^2 = h\nu' + \sqrt{(p_e c)^2 + (m_e c^2)^2}.
\]
Combining this with the momentum relations 
and performing a few algebraic transformations yields the well-known \emph{Compton formula}:
\[
\lambda' - \lambda = \frac{h}{m_e c}\,(1 - \cos\theta).
\]

\begin{MathBox}[The Compton Wavelength]
	The quantity
	\[
	\lambda_C = \frac{h}{m_e c} \approx \SI{2.43e-12}{\meter}
	\]
	is called the \textbf{Compton wavelength} 
	and characterizes the scale at which the effect becomes experimentally observable.
\end{MathBox}

% -------------------------------------------------------------
\subsection{Physical Interpretation}
% -------------------------------------------------------------
The change in wavelength depends solely on the scattering angle \( \theta \) 
and is maximal for backscattering (\( \theta = 180^\circ \)). 
This confirms that photons transfer energy and momentum like particles, 
even though they also propagate as waves.

\begin{DidacticBox}[What the Compton Effect Reveals]
	The Compton effect proves that light cannot be understood purely as a wave. 
	It requires a description in terms of \emph{particles with energy and momentum} — 
	the photons. This marked a decisive extension of classical electrodynamics.
\end{DidacticBox}

